% ============================================================
% DRAFT NOTE (related-work)
% Contains: Positioning against KV-cache eviction and learned retention work,
% output-level quality monitoring, and uncertainty/logit-based prediction.
% Written now: The conceptual distinctions per the current-state hazard
% objective.
% Pending: Verification of bibliography records; an uncertainty-estimation
% and hallucination-detection paragraph once we survey that adjacent line.
% Sources: docs/related_work.md; pivot discussion
% ============================================================
\section{Related Work}

KV-cache methods such as StreamingLLM, H2O, SnapKV, and the implementations
collected in kvpress reduce cache cost by retaining, evicting, or otherwise
compressing state \citep{streamingllm,h2o,snapkv,kvpress}. These methods
motivate the setting, but they do not answer the question studied here: while
decoding under a chosen compressor, how much damage is compression inflicting
on this continuation right now?

Learned retention methods are adjacent but distinct. Lookahead Sparse
Attention predicts which compressed chunks will be relevant in a future
window and fetches them ahead of need \citep{lsa}; it is a selection policy
over what to keep. HERALD forecasts the damage a given compression
configuration is causing, without choosing what to retain. A damage forecast
could eventually drive a serving policy, but this study evaluates the
forecast itself rather than a policy.

Output-level quality monitors and post-hoc failure detectors operate on
emitted text, so they fire only after bad text is user-visible. HERALD's
inputs are statistics of the next-token distribution, which the serving stack
computes at every step anyway, so the forecast is available before the
corresponding text quality change materializes.
\todo{Add a paragraph positioning against logit-based uncertainty estimation
and hallucination-detection work, which uses similar features for a different
target.}

Task-level measures such as IFEval instruction accuracy \citep{ifeval} and
the corpus's catastrophic-failure indicators serve here as secondary
validation that distributional damage corresponds to user-relevant failure,
not as supervision. Gradient-boosted trees \citep{xgboost} appear among our
lagged-feature baselines.
